%% Vol II, Chapter 1 — Kernel Architecture
%% SOURCE: .claude/CLAUDE.md (LithosAnanke / StarKernel section)
%%         docs/working/architecture/03-architecture/OVERVIEW.md
%%         src/starkernel/ (37 files per CLAUDE.md)

\chapter{LithosAnanke Kernel Architecture}
\label{vol2:chap:kernel}

%% TODO(bob): opening paragraph. LithosAnanke ("stone inevitability") is the
%% bare-metal UEFI kernel that boots StarForth directly on hardware.
%% Version 1.0.8, monolithic UEFI PE32+ executable. Cite james:2025:lithos.

\section{Design Philosophy}
\label{vol2:sec:kernel:philosophy}

%% TODO(bob): describe the "stone inevitability" naming rationale and the
%% monolithic UEFI PE32+ design decision (vs. split ELF build).

\section{Boot Sequence}
\label{vol2:sec:kernel:boot}

%% SOURCE: .claude/CLAUDE.md (Boot sequence diagram)
%% TODO(bob): expand the boot sequence from the CLAUDE.md diagram into prose
%% with one paragraph per milestone.

The boot sequence from firmware to FORTH REPL:

\begin{enumerate}
  \item UEFI Firmware → \texttt{uefi\_loader.c} (\texttt{BOOTX64.EFI})
  \item \texttt{ExitBootServices()} — kernel owns hardware; receives
        \texttt{BootInfo\{memory map, ACPI, framebuffer\}}
  \item \texttt{kernel\_main()} — M1 through M7 milestone initialization:
    \begin{enumerate}
      \item[M1] Console init (UART 16550 + framebuffer)
      \item[M2] PMM — physical memory manager, bitmap allocator
      \item[M3] VMM — 4-level x86\_64 paging
      \item[M4] IDT + APIC interrupts
      \item[M5] TSC + HPET + APIC timer (100~Hz heartbeat)
      \item[M6] \texttt{kmalloc} heap (16~MB)
      \item[M7] StarForth VM bootstrap + capsule loading → \texttt{ok} REPL
    \end{enumerate}
\end{enumerate}

\section{Milestone Status}
\label{vol2:sec:kernel:milestones}

%% SOURCE: .claude/CLAUDE.md (Milestone status as of v1.0.8)

\begin{table}[ht]
  \centering
  \caption{LithosAnanke v1.0.8 milestone status}
  \label{tab:milestones}
  \begin{tabular}{lll}
    \toprule
    Milestone & Description & Status \\
    \midrule
    M0--M5 & Boot, console, PMM, VMM, IDT/APIC, timer & Complete \\
    M6 & \texttt{kmalloc} infrastructure & Present (full validation deferred) \\
    M7 & StarForth VM bootstrap + capsule execution pipeline & In progress \\
    \bottomrule
  \end{tabular}
\end{table}

\section{Memory Layout}
\label{vol2:sec:kernel:memory}

%% SOURCE: .claude/CLAUDE.md (Kernel memory layout)

\begin{table}[ht]
  \centering
  \caption{Kernel memory layout}
  \label{tab:kernel-memory}
  \begin{tabular}{ll}
    \toprule
    Region & Description \\
    \midrule
    Kernel heap & 16~MB (\texttt{kmalloc}) \\
    Block RAM LBN 0--991 & 1~MB dedicated RAM blocks \\
    Kernel ramdrive LBN 2048--3071 & 1~MB for capsule loading \\
    LBN 2048 & Entry point for \texttt{init.4th} \\
    \bottomrule
  \end{tabular}
\end{table}

\section{Source Tree}
\label{vol2:sec:kernel:source}

%% SOURCE: .claude/CLAUDE.md (src/starkernel/ listing, 37 files)
%% TODO(bob): reproduce the source-tree listing from CLAUDE.md as a
%% structured description (not a raw ls output).

\section{Multi-Architecture Support}
\label{vol2:sec:kernel:arch}

%% TODO(bob): describe amd64, aarch64, riscv64 build targets and
%% the QEMU/OVMF test strategy.
\begin{lstlisting}[language=bash, caption={StarKernel build targets}]
make -f Makefile.starkernel              # Build for amd64
make -f Makefile.starkernel ARCH=aarch64
make -f Makefile.starkernel ARCH=riscv64
make -f Makefile.starkernel qemu        # Run in QEMU with OVMF
\end{lstlisting}

\section{Experimental Campaign Notes}
\label{vol2:sec:kernel:notes}

%% SCRAP: experiments/campaigns/l8_attractor_map/results_20251209_145907/analysis_output/tables/NOTES
%% SOURCE: docs/working/experiments/campaigns/l8_attractor_map/results_20251209_145907/analysis_output/tables/NOTES.md
%% STATUS: CURRENT
%% FITS: vol1-vm-physics/ch-overview, vol2-kernel/ch-overview, vol3-research/ch-overview
%% EDITORIAL: lifted — prose rewritten to press voice

\section{StarForth, LithosAnanke, and StarshipOS: A Technical Overview}

\textit{Author: Robert A. James. Date: 9 April 2026. Status: Working Document.}

\subsection{StarForth}

\subsubsection{What It Is}

StarForth is a pure C99 FORTH-79 virtual machine with no libc dependencies.
It is self-adaptive: it observes its own execution behavior and continuously
adjusts internal parameters to maintain a conservation invariant called K$\equiv$1.0.
It runs on x86\_64, aarch64, and RISC-V.

\subsubsection{Why FORTH}

FORTH is correct for this work for reasons beyond familiarity. FORTH words
are discrete, classifiable, observable computational units — each word has
a definite stack-effect signature, each word consumes a known quantity of
computational energy, each word is a particle with measurable properties.
This is the foundation of StarForth's self-adaptive behavior.

\subsubsection{The Primitive Taxonomy}

StarForth implements 23 word modules (18 FORTH-79 standard, 5 StarForth
extensions). Every primitive word carries quantum-analog properties derived
from its stack-effect signature:

\begin{description}
  \item[Spin.] Derived from net stack-effect arithmetic. A word consuming
        two values and producing one has spin $-\tfrac{1}{2}$; these values
        fall directly from \texttt{(~n~n~--~n~)} notation, as half-integer
        spin falls from the Dirac equation.
  \item[Charge.] The net stack delta. Charge conservation across a complete
        program means the program is stack-balanced.
  \item[Heat.] Execution frequency and recency. Heat decays over time,
        analogous to radioactive decay.
  \item[Mass.] Bytecode footprint of the word definition.
  \item[Color.] Binding state: white words are fully resolved; colored words
        carry unresolved forward references.
\end{description}

Three of these quantities — spin, charge, and color — behave as genuine
conservation laws across well-formed programs. This is a structural property
of the system discovered empirically and subsequently formalized in
Isabelle/HOL.

\subsubsection{The SSM and James Law}

The SSM (Steady State Machine) is the self-adaptive core of StarForth.
It monitors execution behavior through instrumented subsystems and
continuously adjusts internal parameters to maintain K$\equiv$1.0.

%% PATENT: James Law (K≡1.0) and the SSM adaptive mechanism are
%% patent-adjacent. Do not draft claims here.

\textbf{James Law} ($K \equiv 1.0$) is a conservation invariant: the
normalized total execution energy of the StarForth universe is conserved
at unity across all configurations and workloads. Validation basis:
38{,}400 experimental runs spanning 128 factorial parameter configurations,
CV below 1\%, published on SSRN (abstract 5930134).

The SSM operates through four mechanisms:

\begin{description}
  \item[Rolling Window of Truth (RWOT).] A sliding observation window
        encoding five observables: word identity, prev$\to$word transition,
        instruction pointer position, execution heat, and causal stack depth.
        Updated at heartbeat tick intervals, making it a coarse-grained
        lattice observable.
  \item[Adaptive Heartbeat.] A self-adjusting timer governing when the SSM
        evaluates system state and adjusts parameters. The system's intrinsic
        rhythm.
  \item[Physics Hooks.] Every word execution in the inner interpreter loop
        is bracketed by \texttt{physics\_pre\_execute} and
        \texttt{physics\_post\_execute} calls, updating heat, decay, and
        window measurements.
  \item[Decay.] Execution heat decays between heartbeat ticks, keeping the
        observable window relevant to current rather than historical behavior.
\end{description}

\subsubsection{The Phase Space and Attractor}

Six workload types have been studied experimentally: STABLE, VOLATILE,
TRANSITION, TEMPORAL, DIVERSE, and OMNI. Each produces a distinct trajectory
in phase space. All trajectories converge to the same attractor: K$\equiv$1.0.
Phase-space portrait visualizations confirm attractor convergence across all
workload types on all three target architectures.

\subsubsection{The Manifold Navigation Controller (Theoretical)}

Given a bounded system with a known conservation invariant and predictable
attractor manifold, a controller can:

\begin{enumerate}
  \item Identify the system's current phase-space position using three
        observables: RWOT (position), adaptive heartbeat metadata (velocity),
        and spin-class transition rate (acceleration — requires instrumentation).
  \item Compute a perturbance vector sufficient to transport the system from
        its current manifold to an arbitrary target manifold.
  \item Inject the perturbance — the burn — and release.
  \item Allow the system to coast back to K$\equiv$1.0 under its own restoring
        force.
\end{enumerate}

The controller fires once and releases. The conservation invariant guarantees
the return journey. This orbital-mechanics model scales from the word level
through the VM, OS, FPGA, and prospectively to custom silicon.

---

\subsection{LithosAnanke}

\subsubsection{What It Is}

LithosAnanke is a bare-metal operating system for x86\_64, currently at
milestone M7 (VM integration). It boots via UEFI, manages physical and
virtual memory, handles interrupts via IDT/APIC, maintains a timer, manages
a heap, and hosts the StarForth VM as its primary computational substrate.
Written in C99 with no external dependencies.

\subsubsection{The Capsule System}

The primary organizational unit is the \emph{capsule}: a content-addressed,
64-byte cache-line-aligned descriptor. A capsule's ID is its content hash.
Identity is content: two capsules with identical content have identical IDs.
Modification produces a new capsule with a new ID. Capsules are immutable by
design; mutation produces new capsules.

\subsubsection{Component Architecture}

\begin{description}
  \item[Hera.] The foundational VM and heartbeat. The system's ground state.
  \item[Hermes.] Message entropy and TTL decay. Manages the message-passing
        fabric, monitors message lifecycle, and reaps expired messages. The
        natural home of the perturbance controller extension.
  \item[Artemis.] Persistence and event sourcing. Manages the capsule store,
        maintains the event log, and handles durable state.
\end{description}

Components communicate exclusively through message passing. There is no
shared mutable state between components.

\subsubsection{Authentication and Identity}

LithosAnanke implements a hardware-rooted authentication model with no
usernames and no passwords. Identity is carried on a physical thumb drive
containing two raw block partitions.

Access control is governed by temporal ACL grants that decay via the same
Hermes TTL mechanism that governs message lifetime. The default state of
the permission system is no access. Decay to zero requires no explicit
revocation.

\subsubsection{Normalized Virtual Block Space}

From any VM's perspective, LithosAnanke presents a single contiguous block
address space. Physical origin — local drive, system hardware, or remote
cloud mount — is transparent. Blocks appear as grants are issued; blocks
vanish as grants decay. Security is expressed as address-space geometry.

\subsubsection{The Three UX Tiers}

\begin{description}
  \item[CLI.] The traditional FORTH terminal interface. Full power, zero
        overhead. The REPL is operational.
  \item[GUI.] A wireframe soundstage rendered from rasterized vector
        graphics. Words are physical objects; the user assembles colon
        definitions spatially. The FORTH text is always visible underneath.
  \item[VR.] The full holodeck realization. The user stands inside the phase
        space. Words have mass, spin, heat, and charge as observable
        properties. Manifold navigation is literally navigable space.
\end{description}

All three tiers are windows into identical underlying mechanics.

---

\subsection{StarshipOS}

\subsubsection{The SSPU}

The SSPU (Steady State Processing Unit) is a novel processor architecture
extending the StarForth/SSM work into hardware. It communicates entirely via
a message-passing fabric called Gefyra. There are no traditional buses.
Targeted for initial FPGA implementation on a Digilent Genesys ZU
(Zynq UltraScale+).

%% PATENT: SSPU architecture and Gefyra fabric are patent-adjacent.
%% Provisional patent conversion deadline: December 26, 2026.
%% Do not draft claims here.

\subsubsection{Milestone Status}

\begin{center}
\begin{tabular}{lll}
\toprule
Milestone & Description & Status \\
\midrule
M0 & UEFI boot                        & Complete \\
M1 & Physical memory manager          & Complete \\
M2 & Virtual memory manager           & Complete \\
M3 & IDT/APIC interrupt handling       & Complete \\
M4 & Timer                            & Complete \\
M5 & Heap                             & Complete \\
M6 & --                               & Complete \\
M7 & VM integration                   & Active frontier \\
-- & Manifold Navigation Controller   & Theoretical \\
-- & SSPU FPGA implementation         & Planned (Genesys ZU) \\
\bottomrule
\end{tabular}
\end{center}

\subsubsection{Design Philosophy}

The recurring pattern in this architecture is the elimination of problem
domains rather than their simplification. No scheduler, no traditional memory
manager, no username/password authentication, no logout ceremony, no
coordination layer for distributed state — each elimination is a deliberate
architectural decision. The question at each design point was not ``how do we
implement this?'' but ``do we need this at all?''

\subsubsection{Key Terms}

\begin{center}
\begin{tabular}{lp{8cm}}
\toprule
Term & Definition \\
\midrule
K$\equiv$1.0 & James Law — conservation invariant of normalized execution energy \\
SSM   & Steady State Machine — self-adaptive core of StarForth \\
RWOT  & Rolling Window of Truth — coarse-grained phase-space observable \\
SSPU  & Steady State Processing Unit — novel message-passing processor \\
Gefyra & Message-passing fabric connecting SSPU processing elements \\
Capsule & Content-addressed, 64-byte aligned immutable descriptor \\
Hera  & Foundational VM and heartbeat component \\
Hermes & Message entropy, TTL decay, lifecycle management \\
Artemis & Persistence and event sourcing \\
\bottomrule
\end{tabular}
\end{center}

